%!TEX root = ../../thesis-trier.tex

\todo[inline]{viel auskommentiert}

In \cref{sync_max:thm:sync_max_in_PSPACE} we show that deciding
sync-maximality is in \PSPACE.
Then, in \cref{prop:NL-hard}, we show
that each of the problems to decide complete reachability
and to decide sync-maximality are \NL-hard, even for those
without permutational letters (the proof uses one construction to establish \NL-hardness for both problems).
Then in \cref{sync_max:thm:SI-aut_compl_reach_NL-c},
we derive that deciding complete reachability is \NL-complete
for SI-automata.

In~\cite{BondarV16} the question was posed
whether complete reachability can be decided in polynomial time.
This problem has been later solved by Ferens \& Szyku\l a~\cite{DBLP:conf/icalp/FerensS23}.
Still, the precise computational complexity is unknown.

\begin{theorem}[Ferens \& Szyku\l a~\cite{DBLP:conf/icalp/FerensS23}]
	\label{sync_max:thm:complete_reach_in_P}
	Complete reachabilty can be decided in polynomial time.
\end{theorem}

\begin{theorem} \label{sync_max:thm:sync_max_in_PSPACE}
	Given $\semiautomat{A}{Q}{\Sigma}{\delta}$
	it can be decided in \PSPACE\  whether $\mathcal A$ is sync-maximal.
\end{theorem}
\begin{proof}
	We can assume $|Q| \ge 3$.
	First, guessing some $A \subseteq Q$
	with $|A| \ge 2$
	and testing whether it is reachable,
	which can be done in \PSPACE, see page~\pageref{reach-pspace}.
	So, we can decide in \NPSPACE
	whether some subset of size at least two is not reachable.
	Since \NPSPACE = \PSPACE 
	by Savitch's theorem~\cite{DBLP:journals/jcss/Savitch70}
	the complement of this problem, i.e., whether
	all subsets of size at least two are reachable,
	is in \PSPACE.
	By \cref{sync_max:thm:decision:cd_PSPACE},
	testing whether $\mathcal A$
	is completely distinguishable 
	is in \PSPACE.
	Since both properties characterize
	sync-maximality
	for automata with at least three states
	by \crefitem{sync_max:cor:char_sync_max_aut}{cha_sync_max:item:Q3},
	we obtain the claim.
	
	Alternatively, we can also test whether the automaton
	has the structure as in \cref{sync_maximal:thm:structure}
	in \PSPACE\ by \cref{sync_max:thm:complete_reach_in_P,sync_max:thm:decision:cd_PSPACE}
	since~$S$ (and~$q$, $\Gamma$, and $a \in \Sigma$ in  \crefitem{sync_maximal:thm:structure}{sync_maximal:thm:structure:item:2nd})
	can be computed in polynomial time as is easily seen using \cref{prelim:cor:set_sync_words_in_P}.	
\end{proof} 

\begin{theorem}\label{prop:NL-hard} \label{thm:sync_max_NL-hard}
	The two decision problems
	to decide whether a given automaton is completely reachable
	and whether a given automaton is sync-maximal
	are each \NL-hard even for SI-automata without permutational letters as input.
\end{theorem}
\begin{proof}
	We give a reduction from the problem
	to decide whether a given digraph
	is strongly connected, which is \NL-complete~\cite[Cor. 24]{Jones75}, see also \cref{prelim:thm:complexities}.
	So, let $G = (V, E)$ be a digraph. We can assume $|V| \ge 2$.
	Set $\semiautomat{A}{Q}{\Sigma}{\delta}$
	with $Q = V \times \{0,1\}$
	and an alphabet of simple idempotent letters
	\[
	\Sigma = \{ a_{(p, i), (q, i)} \mid (p,q) \in E, i \in \{0,1\}\} \cup 
	\{ a_{(p,0), (p,1)}, a_{(p,1), (p,0)} \mid p \in V \}
	\]
	where $a_{(p, i), (q, j)}$ maps $(p, i)$ to $(q, j)$
	for distinct $(p,i), (q, j) \in Q$, $p,q \in V$, $i,j \in \{0,1\}$.
	Intuitively, we have two copies of $G$
	inside $\mathcal A$, where an edge $(p,q)$ in $G$ corresponds to the letters $ a_{(p, i), (q, i)}$, $i \in \{0,1\}$, and 
	for every two states from $(p,i)$ and $(p,j)$ with $i \ne j$    
	there exist two simple idempotent letters mapping 
	them onto each other.	
	Next, we prove that the following properties are equivalent, which shows that the construction gives a valid reduction for both decision problems: 
	\begin{enumerate} 
		\item \label{prop:NL-hard:1} $G$ is strongly connected, 
		
		\item \label{prop:NL-hard:2} $\mathcal A$ is strongly connected,
		
		\item \label{prop:NL-hard:3} $\mathcal A$ is completely reachable,
		
		%	\item $\mathcal A$ is completely distinguishable, \todo{falsch, compl dist und nicht SCC so bauen liefert compl dist, z.B. wenn alle auf einen zustand, "stern"-automat}
		
		\item \label{prop:NL-hard:4} $\mathcal A$ is sync-maximal. 
	\end{enumerate}
	
	$\ref{prop:NL-hard:1} \Rightarrow \ref{prop:NL-hard:2}$: This is implied by construction.
	
	$\ref{prop:NL-hard:2} \Rightarrow \ref{prop:NL-hard:3}$: As $\mathcal A$ is an SI-automaton without permutational letters,
	this follows from \cref{cor:SI_CR}. 
	
	$\ref{prop:NL-hard:3} \Rightarrow \ref{prop:NL-hard:4}$: Let $A \subseteq Q$ with distinct $(p_1, i), (p_2, j) \in A$ for $i,j\in \{0,1\}$ and, without loss of generality, $(q,0) \in Q \setminus A$.
	There exist %$u_1, u_2 \in \Sigma^*$ with $u_1, u_2 \in \Sigma \cup \{\varepsilon\}$
	$u_1, u_2 \in \Sigma \cup \{\varepsilon\}$ fixing $(q,0)$
	such that $\delta((p_1, i), u_1) = (p_1, 1)$
	and $\delta((p_2, j), u_2) = (p_2, 1)$.
	As $\mathcal A$ is strongly connected by \cref{prelim:lem:compl_reach_implies_SCC_sync},
	we can pick a shortest word  $u \in \Sigma^*$ 
	mapping $(p_1, 1)$ to $(p_2, 1)$.	
	By construction, this word traces a path 
	using only states in $V \times \{1\}$, which gives
	$\delta((q, 0), u) = (q, 0)$,
	$(q, 0) \notin \delta(A, u)$,
	and $\delta((p_1, 1), u) = \delta((p_2, 1), u) = (p_2, 1)$.
	Then for $v = u_1 u_2 u$
	we have $|\delta(A, v)| < |A|$ and $\delta((q, 0),v) \notin \delta(A, v)$.
	So, by \crefitem{sync_max:thm:char_distinguishable}{dist:item:compress}, 
	$\mathcal A$ is completely distinguishable.
	Hence, by \crefitem{sync_maximal:thm:structure}{sync_maximal:thm:structure:1}
	the automaton is sync-maximal.
	
	$\ref{prop:NL-hard:4} \Rightarrow \ref{prop:NL-hard:1}$:  Assume $G$ is not strongly connected.
	Then there exist distinct $p,q \in V$ such that we have no path from $p$ to $q$.
	By construction, there exists no path from $(p,0)$ to $(q,0)$
	and no path from $(p,1)$ to $(q,1)$.
	By \cref{sync_maximal:thm:structure}
	a sync-maximal automaton has at most one state
	that is not reachable from some other state.
	So, $\mathcal A$ is not sync-maximal.   
	\begin{comment} 
		By construction, (1) implies (2).
		As $\mathcal A$ is an SI-automaton without permutational letters,
		(2) implies (3).
		Next, suppose (3) is true.
		Let $A \subseteq Q$ with distinct $(p_1, i), (p_2, j) \in A$ for $i,j\in \{0,1\}$ and, without loss of generality, $(q,0) \in Q \setminus A$.
		There exist $u_1, u_2 \in \Sigma^*$ with $|u_1|, |u_2| \le 1$
		such that $\delta((p_1, i), u_1) = (p_1, 1)$
		and $\delta((p_2, j), u_2) = (p_2, 1)$.
		As $\mathcal A$ is strongly connected by \cref{prelim:lem:compl_reach_implies_SCC_sync},
		we can pick a shortest word  $u \in \Sigma^*$ 
		mapping $(p_1, 1)$ to $(p_2, 1)$.	
		By construction, this word traces a path 
		using only states in $V \times \{1\}$, which gives
		$\delta((q, 0), u) = (q, 0)$,
		$(q, 0) \notin \delta(A, u)$,
		and $\delta((p_1, 1), u) = \delta((p_2, 1), u) = (p_2, 1)$.
		Then for $v = u_1 u_2 u$
		we have $|\delta(A, v)| < |A|$ and $\delta((q, 0),v) \notin \delta(A, v)$
		and, by \crefitem{sync_max:thm:char_distinguishable}{dist:item:compress}, 
		$\mathcal A$ is completely distinguishable.
		So, by \crefitem{sync_maximal:thm:structure}{sync_maximal:thm:structure:1}
		the automaton is sync-maximal.
		%	Now, assume (4) is true.
		%    By \cref{sync_max:prop:dist_sync}, $\mathcal A$
		%    is synchronizing.
		%    Assume $\mathcal A$ is not strongly connected.
		%    Then there exist distinct $p,q \in V$
		%    such that between $(p,0)$, $(q, 0)$ and $(p, 1)$ and $(q,1)$
		%    there is no path (recall that by construction there is an edge between $(p,0)$ and $(p,1)$),
		%   % By construction a shortest word mapping two states $(p,0)$, $(q, 0)$
		%  %  to the same state induces a path using 
		%    which implies that there exists no word \todo{genauer?}
		%    mapping them to a common word
		%    and so $\mathcal A$ is not synchronizing. \todo{by theorem...}
		%    Hence, $\mathcal A$ is strongly connected
		%    and by \cref{thm:SI_CR} we obtain
		%    that it is completely reachable.
		%    So, by \cref{sync_maximal:thm:structure}, 
		%    it follows that $\mathcal A$ is sync-maximal.
		%    Next, suppose (5) is true.
		%   Then $\mathcal A$ is synchronizing and, similarly as above, we obtain that
		%   $\mathcal A$ must be strongly connected.
		%	By construction $\mathcal A$ has no ``dangling'' state
		%	as in \crefitem{sync_maximal:thm:structure}{sync_maximal:thm:structure:item:2nd}
		%	so that it is completely reachable and so strongly connected by \cref{prelim:lem:compl_reach_implies_SCC_sync}.
		%	By construction, this implies that $G$ 
		%	is strongly connected. \todo{genauer?} 
	\end{comment}
\end{proof}

\begin{theorem}\label{sync_max:thm:SI-aut_compl_reach_NL-c}
	Deciding complete reachability is \NL-complete
	for SI-automata.
\end{theorem}
\begin{proof}
	By \cref{sync_max:cor:SIcr}, for SI-automata,
	complete reachability is equivalent
	to synchronizability and strong connectedness,
	and each property can both be decided in~\NL\ by \cref{prelim:thm:complexities,prelim:complexity:general_aut_sync}.
	By \cref{thm:cd_NL-hard}, deciding complete reachability
	is \NL-hard even for SI-automata as input.
\end{proof}